\chapter{What problem is the paper trying to solve?}\label{ch:story}
\paperloc{Title, abstract, Introduction, and Conclusion.}

\section{Printing a pattern with light}
A semiconductor chip is made by repeatedly creating patterns in thin
layers on a wafer. In optical lithography, light illuminates a patterned
mask, an optical system projects a smaller image onto a light-sensitive
material called photoresist, and chemical processing turns the exposure
into a material pattern. Etching or deposition then transfers or uses that
pattern. The paper concerns the optical exposure part of this sequence.

A \emph{contact hole} is an opening intended to connect device regions or
conducting layers. A \emph{via} is also an interconnection opening, commonly
between conducting levels. For the optical argument, both are small
two-dimensional features whose position and size must be controlled.
A square or rectangle in the mask does not automatically print as an
identical square or rectangle: diffraction rounds corners and mixes light
from nearby features.

The \emph{critical dimension}, abbreviated CD, is a specified feature width.
A rectangle needs at least a horizontal width $\mathrm{CD}_x$ and a vertical
width $\mathrm{CD}_y$. \emph{Pitch} is the centre-to-centre repetition
distance of a periodic pattern. It is not the width of a hole and not the
space between neighbouring edges. For a 14 nm square hole on 38 nm pitch,
the nominal edge-to-edge gap along an axis is $38-14=24\nm$.

The name \emph{5 nm technology node} does not mean that every printed line,
hole, or transistor dimension is 5 nm. It labels a technology generation.
The authors use nominal hole widths of 14 nm and 28 nm. The 2021 IRDS
lithography roadmap explicitly separates node labels from physical
dimensions \cite{irds}. Always identify the layer and feature before
assigning a physical length to a node name.

\section{What makes an image imperfect?}
Even a perfectly manufactured finite optical system has diffraction.
That unavoidable spreading is different from \emph{aberration}, the
departure of an actual optical wavefront from a chosen ideal wavefront.
Aberrations change the phases with which light from different pupil
locations contributes to an image.

The important objects form a hierarchy:
\begin{enumerate}
\item \textbf{Physical hardware:} mirror shapes, positions, coatings,
the mask stack, and illumination optics.
\item \textbf{Optical response:} amplitude, phase, polarization, and the
diffraction orders admitted by the pupil.
\item \textbf{Image:} a continuous distribution of optical intensity.
\item \textbf{Printed geometry:} edges selected by the resist process.
\item \textbf{Performance indicators:} widths, shifts, and changes under
focus and exposure variations.
\end{enumerate}
A neural network operating on the final indicators does not directly
specify a set of manufacturable mirror surfaces. It proposes a wavefront
description that still needs an optical realization.

\section{Why a budget is useful}
Suppose a contact must remain within $0.4\nm$ of its intended vertical
position. Many optical errors contribute to that position. A budget asks
how much of each error can be tolerated while satisfying the performance
requirement. A conventional budget might restrict the total RMS wavefront
error. This paper instead explores combinations of wavefront coefficients
that are acceptable for a particular contact-layer test set.

Write a vector of aberration coefficients as
$\vct c=(c_1,\ldots,c_K)^T$ and the computed imaging indicators as
$\vct y=F(\vct c)$. The symbol $F$ stands for an entire optical and
image-measurement calculation. The forward task calculates $\vct y$
from known $\vct c$. The inverse task finds some $\vct c$ giving acceptable
$\vct y$.

The paper trains an approximate inverse map $G$:
\[
\vct c_{\rm candidate}=G(\vct y_{\rm desired}),\qquad
F(\vct c_{\rm candidate})\approx\vct y_{\rm desired}.
\]
It then sends the proposed coefficients back through the commercial
forward simulator. That final comparison is essential: a network output
is a candidate, and its actual simulated image determines whether it works.

\physicalmeaning{The central proposal is to choose allowable wavefronts using
the imaging needs of a particular layer. The contribution is the
layer-specific inverse modelling and candidate-screening workflow,
not a new law of diffraction.}

\section{Vocabulary you will repeatedly meet}
\begin{longtable}{@{}p{.23\textwidth}p{.71\textwidth}@{}}
\toprule Term & Meaning in this study\\ \midrule\endhead
EUV & Extreme ultraviolet light; here the wavelength is 13.5 nm.\\
DUV & Deep ultraviolet; the Introduction compares EUV with 193 nm light.\\
Projector / objective & The optical system forming the mask image on
the wafer; in EUV it uses reflective optics.\\
Pupil & A representation of the aperture selecting ray directions or
spatial-frequency components.\\
Wavefront & A surface of equal optical phase at a specified time.\\
OPD / WFE & Optical path difference / wavefront error relative to a
specified ideal reference, usually expressed as a length.\\
Zernike coefficient & The amplitude of one specified polynomial mode
in a pupil-wavefront expansion.\\
PS & Pattern shift: displacement of a specified printed feature centre.\\
PVB & Process variation band: the spread of printed contours when
specified process conditions change.\\
NILS & Normalized image log slope: relative intensity steepness at an edge.\\
BPNN & Backpropagation neural network, a trainable nonlinear function.\\
OPC & Optical proximity correction: modification of the mask layout to
improve the eventual printed pattern.\\
\bottomrule
\end{longtable}
\checkpoint{If two patterns have the same CD, must they have the same
pitch, centre position, or process robustness? No. These are separate
properties; this is why the paper uses several indicators for many patterns.}

\chapter{The mathematics we actually need}\label{ch:math}
\section{Units and dimensionless ratios}
One nanometre is $10^{-9}$ metres; one micrometre is $10^{-6}$ metres.
An angle in radians is the ratio of arc length to radius and is dimensionless.
An exponential must have a dimensionless argument.
For example, $2\pi W/\lambda$ is legitimate because $W$ and $\lambda$
are both lengths.

A \emph{wave} of optical path error means one wavelength of OPD.
A milliwave, written here as $\mathrm{m}\lambda$, is $10^{-3}$ waves:
\begin{equation}
W_{\rm nm}=\lambda_{\rm nm}\frac{c_{\mathrm{m}\lambda}}{1000}.
\label{eq:milliwave}
\end{equation}
The paper's MMW is used in this milliwave sense. It is not a millimetre.
At 13.5 nm, $1\,\mathrm{m}\lambda=0.0135\nm$ and
$20\,\mathrm{m}\lambda=0.270\nm$.

\section{Complex numbers are a language for phase}
The imaginary unit obeys $\ii^2=-1$. A complex number $z=a+\ii b$ has
conjugate $z^*=a-\ii b$, modulus $|z|=\sqrt{a^2+b^2}$, and squared modulus
$zz^*=a^2+b^2$. Euler's formula is
\begin{equation}
\ee^{\ii\phi}=\cos\phi+\ii\sin\phi.
\end{equation}
Multiplication by $\ee^{\ii\phi}$ rotates a complex number by $\phi$.
It changes phase without changing magnitude.

Represent a real harmonic electric field by
\begin{equation}
E(\vct r,t)=\Real\{U(\vct r)\ee^{-\ii\omega t}\},\qquad
U=A\ee^{\ii\phi}.
\label{eq:phasor}
\end{equation}
The observable field is real. The complex quantity $U$ stores its
amplitude and phase compactly. Time differentiation becomes multiplication
by $-\ii\omega$, a major simplification.

For two scalar fields $U_1$ and $U_2$, expand the intensity:
\begin{align}
|U_1+U_2|^2
&=(U_1+U_2)(U_1^*+U_2^*)\\
&=|U_1|^2+|U_2|^2+U_1U_2^*+U_1^*U_2\\
&=A_1^2+A_2^2+2A_1A_2\cos(\phi_1-\phi_2).
\label{eq:interference}
\end{align}
The last term is interference. Equal amplitudes in phase give intensity
$4A^2$; equal amplitudes out of phase by $\pi$ give zero. The separate
intensities were $A^2$ in both cases. This is why changing phase can
change an image without changing the amount of light admitted by the pupil.

\section{Derivatives, gradients, and Taylor expansion}
A derivative measures a local change. For small displacement $\delta x$,
\[
I(x+\delta x)=I(x)+I'(x)\delta x+O(\delta x^2).
\]
Here $O(\delta x^2)$ denotes terms whose leading size is proportional
to the square of the small displacement. A partial derivative varies
one argument while holding the others fixed:
\[
\nabla_\perp W=\left(\frac{\partial W}{\partial x},
\frac{\partial W}{\partial y}\right).
\]
For $W=A(x^2+y^2)^2$,
$\nabla_\perp W=4A(x^2+y^2)(x,y)$.
If $W$ is a length, this gradient is dimensionless; it will be related
to an angular error.

We frequently use
\begin{align}
\sqrt{1+\epsilon}&=1+\frac{\epsilon}{2}
-\frac{\epsilon^2}{8}+O(\epsilon^3),\\
\ee^{\ii\phi}&=1+\ii\phi-\frac{\phi^2}{2}+O(\phi^3),\\
\ln(1+\epsilon)&=\epsilon+O(\epsilon^2).
\end{align}
These follow by evaluating successive derivatives at zero. They are
approximations for small dimensionless arguments. Calling a formula
paraxial or small-aberration specifies the expansion that has been made;
it does not remove the need to check the argument's size.

\worked{Changing to normalized pupil coordinates}
Let $x_p=a u$, $y_p=a v$, where $a$ is the physical pupil radius,
and $\widetilde W(u,v)=W(au,av)$. The chain rule gives
\begin{equation}
\frac{\partial W}{\partial x_p}
=\frac{1}{a}\frac{\partial\widetilde W}{\partial u}.
\label{eq:chain}
\end{equation}
For $\widetilde W=B u^2$, the physical slope is $2Bu/a$,
not $2Bu$. Omitting $1/a$ changes both dimensions and physics.

\section{Integrals over a disk}
Use normalized polar coordinates $u=\rho\cos\theta$,
$v=\rho\sin\theta$, $0\leq\rho\leq1$.
The small area is $\dd u\,\dd v=\rho\,\dd\rho\,\dd\theta$.
The uniformly weighted disk average is
\begin{equation}
\avg{f}=\frac{1}{\pi}\int_0^{2\pi}\int_0^1
f(\rho,\theta)\rho\,\dd\rho\,\dd\theta.
\label{eq:diskaverage}
\end{equation}
The factor $\rho$ accounts for the growing circumference of rings.
Sampling equally spaced radii with equal weights is not uniform area sampling.
For $q>-2$,
\begin{equation}
\avg{\rho^q}=2\int_0^1\rho^{q+1}\dd\rho=\frac{2}{q+2}.
\label{eq:diskpowers}
\end{equation}
Thus $\avg{\rho^2}=1/2$, $\avg{\rho^4}=1/3$, and
$\avg{\rho^8}=1/5$. Also $\avg{\rho^2\cos^2\theta}=1/4$
because the angular average of $\cos^2\theta$ is $1/2$.

The mean-subtracted RMS of a real wavefront is
\begin{equation}
\sigma_W=\sqrt{\avg{(W-\avg W)^2}}
=\sqrt{\avg{W^2}-\avg W^2}.
\label{eq:variance}
\end{equation}
It is a length. The mean $\avg W$ is called piston.
For $W=B\rho^2$, $\sigma_W=|B|/\sqrt{12}$.

\section{Vectors, matrices, and least squares}
A vector is an ordered list. A matrix applies weighted sums to that list.
If $A$ has $M$ rows and $K$ columns, $A\vct c$ has $M$ entries:
\[
(A\vct c)_i=\sum_{j=1}^K A_{ij}c_j.
\]
The transpose $A^T$ interchanges rows and columns. The Euclidean norm
obeys $\norm{\vct c}^2=\vct c^T\vct c=\sum_jc_j^2$.
To fit data $\vct b$ with $A\vct c$, minimize
$L=\tfrac12\norm{A\vct c-\vct b}^2$.
Expansion and differentiation give
$\nabla_{\vct c}L=A^T(A\vct c-\vct b)$.
At a stationary point,
\begin{equation}
A^TA\,\vct c=A^T\vct b.
\label{eq:normal}
\end{equation}
These are the normal equations. A unique solution requires independent
columns. If a nonzero $\vct v$ satisfies $A\vct v=0$, then $\vct c$
and $\vct c+\vct v$ predict the same data. Such a vector is in the
\emph{null space}. This idea becomes central to the paper's claim that
different aberration combinations can have similar imaging performance.

\section{Averages, random errors, and correlations}
For samples $x_1,\ldots,x_N$, the mean is $\bar x=N^{-1}\sum_i x_i$.
An expectation $\E[x]$ is the corresponding probability-weighted average.
Variance is $\E[(x-\E[x])^2]$ and
\[
\Cov(x,y)=\E[(x-\E[x])(y-\E[y])].
\]
Expanding the square proves
\begin{equation}
\Var(x+y)=\Var(x)+\Var(y)+2\Cov(x,y).
\label{eq:covsum}
\end{equation}
The root-sum-square rule follows only when covariance vanishes.
Two fabrication errors that always move together do not behave like
independent errors. Fixed errors do not become random simply because
their signs are unknown.

\section{Fourier analysis: a shape as a sum of ripples}
A periodic function of period $p$ can be represented as
\begin{equation}
t(x)=\sum_{m=-\infty}^{\infty}a_m\ee^{\ii2\pi mx/p},\qquad
a_m=\frac1p\int_{-p/2}^{p/2}t(x)\ee^{-\ii2\pi mx/p}\dd x.
\label{eq:fourierseries}
\end{equation}
Substitute the series into the integral. The integral of
$\ee^{\ii2\pi(m-n)x/p}$ over one period is zero unless $m=n$, when it
equals $p$. The integral selects one ripple.

For nonperiodic two-dimensional fields we use
\begin{align}
\widetilde U(\vct f)&=\int_{\R^2}U(\vct x)
\ee^{-\ii2\pi\vct f\cdot\vct x}\dd^2x,\\
U(\vct x)&=\int_{\R^2}\widetilde U(\vct f)
\ee^{\ii2\pi\vct f\cdot\vct x}\dd^2f.
\label{eq:fourierpair}
\end{align}
Spatial frequency $\vct f=(f_x,f_y)$ has units of inverse length.
Fine detail needs higher spatial frequencies. A finite pupil removes
some frequencies, so an image cannot reproduce arbitrary fine detail.

Convolution, $(h*t)(x)=\int h(x-x')t(x')\dd x'$, means summing shifted
responses to all object points. Substitution of the Fourier pair shows
$\FT[h*t]=\widetilde h\,\widetilde t$: convolution in position becomes
multiplication in frequency. This is the algebraic foundation of Fourier
optics, not a separate empirical imaging rule.

\chapter{Light, phase, rays, and optical path}\label{ch:waves}
\reading{Saleh--Teich \S\S2.1--2.2, printed pp.~40--47 (PDF pp.~62--69),
and \S2.5, p.~58 (PDF p.~80). Born--Wolf \S1.3, p.~14 (PDF p.~43),
and \S3.1, p.~109 (PDF p.~138).}

\section{From Maxwell's equations to a wave}
In a homogeneous, source-free, linear, isotropic, nonmagnetic medium,
Maxwell's equations include
\[
\nabla\times\vct E=-\mu_0\frac{\partial\vct H}{\partial t},\quad
\nabla\times\vct H=\epsilon\frac{\partial\vct E}{\partial t},\quad
\nabla\cdot\vct E=0.
\]
Take the curl of the first equation and substitute the second.
Using $\nabla\times(\nabla\times\vct E)
=\nabla(\nabla\cdot\vct E)-\nabla^2\vct E$ gives
\begin{equation}
\nabla^2\vct E-\mu_0\epsilon\frac{\partial^2\vct E}{\partial t^2}=0.
\end{equation}
The wave speed is $v=1/\sqrt{\mu_0\epsilon}=c/n$.
For one frequency and a scalar component under a justified approximation,
\begin{equation}
\nabla^2 U+n^2k_0^2U=0,\qquad
k_0=\frac{2\pi}{\lambda_0}=\frac{\omega}{c}.
\label{eq:helmholtz}
\end{equation}
This is the Helmholtz equation. A solution
$U=A\ee^{\ii\vct k\cdot\vct r}$ is a plane wave if
$|\vct k|=nk_0$. Constant-phase surfaces are planes perpendicular to
$\vct k$. given that $|\vct k|=nk_0$, $\lambda_{medium}  = \frac{2\pi}{|\vct k|} = \frac{\lambda}{n}$

\subsection*{From the Poynting vector to optical intensity}
We now derive the intensity formula rather than introduce it as a
separate optical rule. The result for a travelling plane wave in a
lossless, nonmagnetic dielectric is
\[
I=\frac{n\epsilon_0c}{2}|\vct E_0|^2.
\]
Here $\vct E_0$ is an electric-field phasor, measured in volts per metre,
and its components are defined using peak rather than RMS amplitudes.
The factor $1/2$ comes from averaging the real oscillating fields.
The factor $n\epsilon_0c$ comes from the relation between the electric
and magnetic fields. We will obtain these factors separately.

\subsubsection*{1. The Poynting vector measures energy flow}
For real, instantaneous fields in SI units, the Poynting vector is
\begin{equation}
\vct S(\vct r,t)=\vct E(\vct r,t)\times\vct H(\vct r,t)
=\frac{1}{\mu_0}\vct E(\vct r,t)\times\vct B(\vct r,t),
\label{eq:poynting-instant}
\end{equation}
where $\vct B=\mu_0\vct H$ in the nonmagnetic medium considered here.
The cross product gives a direction of energy flow. Its units are
\[
[\vct S]=\frac{\mathrm V}{\mathrm m}
          \frac{\mathrm A}{\mathrm m}
=\frac{\mathrm W}{\mathrm m^2},
\]
so it is power per area. For a surface element $\dd A$ with unit normal
$\widehat{\vct n}_A$, the signed power crossing it is
$\vct S\cdot\widehat{\vct n}_A\,\dd A$.

Why does this particular cross product describe energy?
For a lossless, nondispersive dielectric with constant $\epsilon$,
use the vector identity
\[
\nabla\cdot(\vct E\times\vct H)
=\vct H\cdot(\nabla\times\vct E)
-\vct E\cdot(\nabla\times\vct H).
\]
Substituting the source-free Maxwell equations above gives
\begin{align}
\nabla\cdot\vct S
&=-\mu_0\vct H\cdot\frac{\partial\vct H}{\partial t}
  -\epsilon\vct E\cdot\frac{\partial\vct E}{\partial t}
\notag\\
&=-\frac{\partial}{\partial t}
  \left(\frac{\epsilon|\vct E|^2}{2}
             +\frac{\mu_0|\vct H|^2}{2}\right).
\end{align}
The expression in parentheses is the electromagnetic energy density
$u_{\rm em}$, in joules per cubic metre. Thus
\begin{equation}
\frac{\partial u_{\rm em}}{\partial t}+\nabla\cdot\vct S=0,
\qquad
\frac{\dd}{\dd t}\int_Vu_{\rm em}\,\dd V
=-\oint_{\partial V}\vct S\cdot\widehat{\vct n}_A\,\dd A.
\label{eq:poynting-conservation}
\end{equation}
The second equality follows from the divergence theorem: outward energy
flow through a boundary reduces the energy stored inside it.
This identifies $\vct S$ as an energy-flux density.
The simple instantaneous formula for $u_{\rm em}$ assumes no material
dispersion; the monochromatic flux calculation below only needs the
real, lossless material parameters at the frequency of interest.

\subsubsection*{2. Average the real fields, not just their phasors}
An optical detector or resist exposure averages over many optical cycles.
Let $T_{\rm opt}=2\pi/\omega$ and define a one-cycle average at a fixed point:
\[
\avg{\vct S}_t
=\frac{1}{T_{\rm opt}}\int_{t_0}^{t_0+T_{\rm opt}}\vct S(t)\,\dd t.
\]
For an exactly monochromatic field this equals an average over any
integer number of cycles. Write the real fields in terms of their
complex spatial phasors:
\begin{align*}
\vct E(\vct r,t)
&=\Real\{\widetilde{\vct E}(\vct r)\ee^{-\ii\omega t}\}
=\frac12\left(
\widetilde{\vct E}\ee^{-\ii\omega t}
+\widetilde{\vct E}^{\,*}\ee^{+\ii\omega t}\right),\\
\vct H(\vct r,t)
&=\Real\{\widetilde{\vct H}(\vct r)\ee^{-\ii\omega t}\}
=\frac12\left(
\widetilde{\vct H}\ee^{-\ii\omega t}
+\widetilde{\vct H}^{\,*}\ee^{+\ii\omega t}\right).
\end{align*}
Suppressing the common position argument, their cross product is
\begin{align*}
\vct E\times\vct H
=\frac14\Bigl[&
(\widetilde{\vct E}\times\widetilde{\vct H})\ee^{-2\ii\omega t}
+\widetilde{\vct E}\times\widetilde{\vct H}^{\,*}\\
&+\widetilde{\vct E}^{\,*}\times\widetilde{\vct H}
+(\widetilde{\vct E}^{\,*}\times\widetilde{\vct H}^{\,*})
\ee^{+2\ii\omega t}\Bigr].
\end{align*}
The first and last terms average to zero because
$\avg{\ee^{\pm2\ii\omega t}}_t=0$.
The two remaining terms are complex conjugates of one another.
Since $z+z^*=2\Real z$, applied component by component,
\begin{equation}
\boxed{\avg{\vct S}_t
=\frac12\Real\!\left(
\widetilde{\vct E}\times\widetilde{\vct H}^{\,*}\right).}
\label{eq:poynting-average}
\end{equation}
The conjugation and the real part are essential: an unmodified cross
product of complex phasors is not the measured average power flow.
This expression also applies to monochromatic fields that are not plane
waves; the following reduction to electric-field magnitude alone uses
the plane-wave assumptions.

\subsubsection*{3. Maxwell's equations relate the two field amplitudes}
For a travelling plane wave, write
\[
\widetilde{\vct E}(\vct r)
=\vct E_0\ee^{\ii\vct k\cdot\vct r},\qquad
\widetilde{\vct H}(\vct r)
=\vct H_0\ee^{\ii\vct k\cdot\vct r},
\qquad
\vct k=k\widehat{\vct s},
\]
where $\widehat{\vct s}$ is a real unit vector in the propagation direction
and $k=n\omega/c$.
For the time convention $\ee^{-\ii\omega t}$, Faraday's law becomes
\[
\nabla\times\widetilde{\vct E}
=\ii\omega\mu_0\widetilde{\vct H}.
\]
Taking the curl of the plane-wave exponential gives
$\nabla\times\widetilde{\vct E}
=\ii\vct k\times\widetilde{\vct E}$.
Cancel the common exponential and $\ii$ to obtain
\begin{equation}
\vct H_0
=\frac{\vct k\times\vct E_0}{\omega\mu_0}
=\frac{1}{\eta}\widehat{\vct s}\times\vct E_0,
\qquad
\eta=\sqrt{\frac{\mu_0}{\epsilon}}
=\frac{\eta_0}{n},\qquad
\eta_0=\sqrt{\frac{\mu_0}{\epsilon_0}}.
\label{eq:plane-wave-impedance}
\end{equation}
We used $\epsilon=n^2\epsilon_0$ and
$c=1/\sqrt{\mu_0\epsilon_0}$.
The quantity $\eta$ is the wave impedance: it sets the ratio of electric
to magnetic field amplitude and has units of ohms.
Also, $\nabla\cdot\widetilde{\vct E}=0$ gives
$\widehat{\vct s}\cdot\vct E_0=0$, so the electric field is transverse.
The magnetic field is transverse as well, and its direction follows from
the cross product in Eq.~\eqref{eq:plane-wave-impedance}.

\subsubsection*{4. Obtain the intensity and its direction}
Because $\vct k$ is real, the spatial factors
$\ee^{\ii\vct k\cdot\vct r}$ and
$\ee^{-\ii\vct k\cdot\vct r}$ cancel in
Eq.~\eqref{eq:poynting-average}. Because $\eta$ is real,
\begin{align*}
\avg{\vct S}_t
&=\frac{1}{2\eta}
\Real\!\left[\vct E_0\times
             (\widehat{\vct s}\times\vct E_0^*)\right]\\
&=\frac{1}{2\eta}
\Real\!\left[
\widehat{\vct s}(\vct E_0\cdot\vct E_0^*)
-\vct E_0^*(\vct E_0\cdot\widehat{\vct s})\right]\\
&=\frac{|\vct E_0|^2}{2\eta}\widehat{\vct s}.
\end{align*}
The second line is the vector triple-product identity.
Transversality removes its second term, while
$|\vct E_0|^2=\vct E_0\cdot\vct E_0^*
=|E_{0x}|^2+|E_{0y}|^2+|E_{0z}|^2$ is real and nonnegative.
Define the plane-wave intensity as the average power per area
perpendicular to propagation:
\begin{equation}
\boxed{
I=\avg{\vct S}_t\cdot\widehat{\vct s}
=\frac{|\vct E_0|^2}{2\eta}
=\frac{n\epsilon_0c}{2}|\vct E_0|^2.}
\label{eq:plane-wave-intensity}
\end{equation}
For the last equality,
$1/\eta=n/\eta_0=n\sqrt{\epsilon_0/\mu_0}=n\epsilon_0c$.
The result holds for linear, elliptical, and circular polarization in
this medium; it uses the sum of the squared component phasor magnitudes.

If a detector surface normal makes angle $\alpha$ with
$\widehat{\vct s}$, its incident power per detector area is
$\avg{\vct S}_t\cdot\widehat{\vct n}_A=I\cos\alpha$, taking the normal
toward the direction of energy flow. Thus one must specify the area
orientation when converting flux into detected power.

\subsubsection*{5. Check the factor of two with a real cosine wave}
Take a linearly polarized wave travelling along $+z$:
\[
\vct E=\widehat{\vct x}E_{\rm pk}\cos(kz-\omega t),\qquad
\vct H=\widehat{\vct y}\frac{E_{\rm pk}}{\eta}
\cos(kz-\omega t).
\]
Then
\[
\vct S=\widehat{\vct z}\frac{E_{\rm pk}^2}{\eta}
\cos^2(kz-\omega t),\qquad
\avg{\cos^2(kz-\omega t)}_t=\frac12,
\]
which gives $I=E_{\rm pk}^2/(2\eta)$ directly.
If the field is instead specified by
$E_{\rm rms}=E_{\rm pk}/\sqrt2$, the same physical result is
$I=E_{\rm rms}^2/\eta$.
The two formulas use different amplitude conventions, so the factor
$1/2$ must not be inserted twice.
As a units-and-scale check, a vacuum plane wave with
$E_{\rm pk}=1\,\mathrm{V/m}$ has
$I=1/(2\eta_0)\simeq1.33\times10^{-3}\,\mathrm{W/m^2}$,
using $\eta_0\simeq376.73\,\Omega$.

\subsubsection*{6. Why scalar optics can write \texorpdfstring{$I=|U|^2$}{I = |U| squared}}
For a fixed unit polarization vector $\widehat{\vct e}$, write
$\widetilde{\vct E}=\widehat{\vct e}\,\mathcal E$ with
$\widehat{\vct e}\cdot\widehat{\vct e}^{\,*}=1$.
Here $\mathcal E$ is the scalar electric-field phasor in volts per metre.
Under the travelling-wave relation just derived,
\[
I=\frac{n\epsilon_0c}{2}|\mathcal E|^2.
\]
Define a rescaled scalar amplitude
\begin{equation}
U=\sqrt{\frac{n\epsilon_0c}{2}}\,\mathcal E
\quad\Longrightarrow\quad I=|U|^2.
\label{eq:intensity-normalized-field}
\end{equation}
The constant is absorbed into the \emph{amplitude through its square root}.
This $U$ has units $\mathrm{W}^{1/2}/\mathrm m$ if $I$ is in
$\mathrm{W/m^2}$; it is no longer the numerical electric field in
$\mathrm{V/m}$.
Multiplying by a spatially constant factor leaves the homogeneous
Helmholtz equation unchanged. In a scalar, paraxial imaging model with
a common polarization and nearly common propagation direction, this
normalization is why we may add the scalar fields and then take
$|U|^2$ to obtain intensity.

One may also divide intensity by a fixed reference intensity and use
dimensionless fields. Neither rescaling changes the interference or
edge-position argument, provided the threshold and dose use the same
normalization. For absorbing media, standing waves, or fields with
substantially different propagation directions, return to the full
flux expression in Eq.~\eqref{eq:poynting-average}; the simple
real-$n$ plane-wave prefactor is not a universal replacement for it.

\section{Optical path is not just distance}
Along a geometrical trajectory of arc length $\ell$,
\begin{equation}
\mathcal L=\int n(\vct r)\,\dd\ell.
\label{eq:opl}
\end{equation}
A uniform slab of thickness $t$ and index $n$ has optical path $nt$.
Phase accumulated along a ray is $k_0\mathcal L$ in the
$\ee^{-\ii\omega t}$ convention. Comparing actual and reference optical
paths gives OPD $W$ and phase error
\begin{equation}
\phi_W=k_0 W=\frac{2\pi W}{\lambda_0}.
\label{eq:phaseOPD}
\end{equation}
A reference must be specified: usually an ideal spherical wave converging
to a chosen image point. Moving that reference point changes the
reported tilt and defocus, even when the physical system is unchanged.

\worked{The wavelength comparison in the Introduction}
For fixed physical RMS OPD of 1.5 nm,
\[
\frac{1.5}{193}=0.0077720\ {\rm waves}
=7.7720\,\mathrm{m}\lambda,\qquad
\frac{1.5}{13.5}=0.111111\ {\rm waves}
=111.111\,\mathrm{m}\lambda.
\]
The ratio is $193/13.5=14.2963$. The length error has not grown.
Its size measured in wavelengths, and therefore its phase error, has grown.
This is the precise meaning of the paper's inverse-wavelength statement.

\section{How rays emerge from wave theory}
Write $U=A(\vct r)\ee^{\ii k_0S(\vct r)}$, where $S$ has length units.
Substitute into the Helmholtz equation:
\[
\nabla^2 A+2\ii k_0\nabla A\cdot\nabla S
+\ii k_0A\nabla^2S-k_0^2A|\nabla S|^2+n^2k_0^2A=0.
\]
When amplitude and index vary slowly on the wavelength scale, the leading
real terms require
\begin{equation}
|\nabla S|^2=n^2.
\label{eq:eikonal}
\end{equation}
This is the eikonal equation. The unit ray direction is
$\vct d=\nabla S/n$. Rays are normal to wavefronts.
The next terms give amplitude transport; ray geometry alone does not
determine interference.

For nearby actual and reference eikonals, $S=S_0+W$, the first-order
transverse direction-cosine difference is
\begin{equation}
\delta\vct d_\perp\simeq\frac{1}{n}\nabla_\perp W.
\label{eq:raygradient}
\end{equation}
Over a paraxial propagation distance $L$, the corresponding displacement
is approximately $L\nabla_\perp W/n$. This is a local geometrical statement.
A finite-aperture lithographic image still requires field summation and
source averaging.

\careful{A coefficient's sign depends on the OPD definition, time convention,
pupil axes, and image inversion. This book uses pupil phase
$\exp(+\ii2\pi W/\lambda)$ and Eq.~\eqref{eq:fourierpair}.
In transfer-function formulas the normalized pupil labels spatial
frequency, $\vct u=\vct f/f_c$; it need not have the orientation of the
physical exit-pupil coordinates in a ray drawing. Image-shift signs below
follow the declared transfer-function convention.}

\section{What a surface error does at a mirror}
For a mirror displaced a small normal distance $h$, the incident and
reflected segments each change optical path by approximately
$h\cos\alpha$, with $\alpha$ measured from the surface normal.
In vacuum and for fixed local ray geometry,
\begin{equation}
|\delta W|\simeq2|h|\cos\alpha.
\label{eq:mirror}
\end{equation}
The sign depends on positive surface displacement. At normal incidence,
a 0.1 nm height error gives about 0.2 nm OPD. Wavefront tolerance and
surface-height tolerance must not be interchanged.

This formula excludes coating phase, changed intersections, multiple-mirror
coordinate mappings, and polarization. Those belong in a full sensitivity
model. A lens-thickness perturbation in air often gives approximately
$(n-1)\delta t$ OPD rather than the mirror factor of two.

\section{Polarization and the limits of a scalar field}
The electric field is a vector perpendicular to the propagation direction
for a plane wave in an isotropic medium. The interference cross term is
$2\Real(\vct U_1\cdot\vct U_2^*)$. Orthogonal polarizations therefore have
no such cross term in total intensity. Reflection gives different
amplitudes and phases to fields perpendicular and parallel to the
incidence plane.

Scalar optics is useful when these vector effects can be separated or
neglected. At an EUV mask with finite-height absorbing features,
wavelength-scale layers, oblique illumination, and polarization-dependent
scattering, a full electromagnetic calculation may be necessary.
Rigorous simulation refers to a more complete numerical treatment of the
specified model; it does not eliminate missing inputs or manufacturing
and resist uncertainty.
